% TOP_PROBLEM: {"id":30007655,"problem_number":"LOCAL-30007655","title":"Microcredit heterogeneity","category_id":21,"category":{"id":21,"name":"economics","display_name":"Economics","description":"Open economic theory statements, published research agendas, and proposed scoped specifications; question provenance and historical priority are qualified per record.","slug":"economics","order_index":21,"created_at":"2026-10-08T00:00:00Z"},"status":"provenance_pending","external_url":null,"legacy_ids":[],"published":false,"statement":"In the explicitly specified setting: Which borrowers and credit-contract designs generate substantial durable business growth rather than modest consumption smoothing, and can those gains scale? The mathematical statement above fixes the target, assumptions and scope of this catalogue formulation.","economics":{"original_id":144,"field":"Development and poverty","kind":"design","formalization_basis":"proposed_specification","source_status":"not_verified","priority_status":"not_established","priority_note":"No qualifying problem statement was directly verified, so historical priority cannot be assigned.","earliest_verified_reference":null,"current_reference":{"authors":["Abhijit Banerjee","Emily Breza","Esther Duflo","Cynthia Kinnan"],"title":"Can Microfinance Unlock a Poverty Trap for Some Entrepreneurs?","year":2019,"url":"https://www.nber.org/papers/w26346","doi":"10.3386/w26346","locator":"Abstract, July 2024 revision","evidence_summary":"The paper answers its local entrepreneurship question with heterogeneous long-run effects; the broader proposed contract-design question is not a verified open statement."},"formalization_note":"This is a proposed scoped research specification. No primary passage explicitly stating this entire remaining question as open was verified. The supporting literature may answer important subquestions; this entry must not be advertised as a verified formal open problem."},"statusQualification":"Provenance pending: an explicit published statement of this exact open question was not verified. This is a catalogue-authored candidate specification, not a certified open conjecture. This is a proposed scoped research specification. No primary passage explicitly stating this entire remaining question as open was verified. The supporting literature may answer important subquestions; this entry must not be advertised as a verified formal open problem.","statusReviewedAt":"2026-10-08"}
% FIRST_POSTED: 2026-10-08
% ENTRY_KIND: statement_only
% !TeX program = lualatex
\documentclass[11pt]{article}
\usepackage{fontspec}
\usepackage[a4paper,margin=25mm]{geometry}
\usepackage{amsmath,amssymb,mathtools}
\usepackage{xurl}
\usepackage[hidelinks]{hyperref}
\newcommand{\sref}[2]{\hyperlink{source:#1:#2}{[S#2]}}
\setlength{\emergencystretch}{3em}
\begin{document}
% problemId:30007655
\section*{Microcredit heterogeneity}\label{30007655}
\subsection{Definitions and mathematical statement}
Fix a population of potential small-business borrowers and a lender with a declared budget. Contract \(p\) specifies credit access, loan size, interest, repayment timing, collateral, and insurance. Let \(D_i(p)\) indicate borrowing, \(\Pi_{iT}(p)\) business profits at horizon \(T\), and \(C_{iT}(p)\) household consumption. Define intention-to-treat effects \(\tau_\Pi(p,x)=E[\Pi_{iT}(p)-\Pi_{iT}(0)\mid X_i=x]\), where \(X_i\) records baseline business activity, assets, and risk. Measure net borrower welfare, lender losses, and default alongside profits. The research task is to identify which observable borrower groups and contract features generate sustained productive gains and which mainly smooth consumption. Offers can be randomized, but borrower-only effects require additional compliance assumptions; conditioning on take-up changes the population. Distinguish learning, indivisible investment, risk reduction, and selection through appropriate variations. Estimate whether scalable selection rules retain gains when repayment incentives and product-market competition respond. Some microcredit effects and heterogeneous responses have already been documented. This is a proposed remaining design and transportability question; the original supporting results are not proof that a precise universal conjecture or first problem statement exists.

\par\textbf{Formulation and scope.} This is a proposed scoped research specification. No primary passage explicitly stating this entire remaining question as open was verified. The supporting literature may answer important subquestions; this entry must not be advertised as a verified formal open problem.

\par\textbf{Open-question provenance.} An explicit published open-question statement was not verified. Sources below are supporting literature and must not be treated as a first listing.

\par\textbf{Historical priority.} No qualifying problem statement was directly verified, so historical priority cannot be assigned.

\subsection{Short English statement}
In the explicitly specified setting: Which borrowers and credit-contract designs generate substantial durable business growth rather than modest consumption smoothing, and can those gains scale? The mathematical statement above fixes the target, assumptions and scope of this catalogue formulation.

\subsection{Sources}
\begin{itemize}\raggedright
\item[\textnormal{[S1]}]\hypertarget{source:30007655:1}{} Abhijit Banerjee, Emily Breza, Esther Duflo, Cynthia Kinnan. 2019. Can Microfinance Unlock a Poverty Trap for Some Entrepreneurs?. Abstract, July 2024 revision.
\url{https://www.nber.org/papers/w26346}.
DOI: 10.3386/w26346.
{\small\textit{Supporting or current literature; see evidence qualification. The paper answers its local entrepreneurship question with heterogeneous long-run effects; the broader proposed contract-design question is not a verified open statement.}}
\end{itemize}
\end{document}
